\section{Gradient stability and polynomial time sampling}
\label{sec:gradstab}

\subsubsection{The computational question} The posterior of a nonlinear partial differential equation inverse problem is a density on $\R^D$ with $D$ large, and it is not log concave, because the forward map is nonlinear and enters the log likelihood through a square. A Markov chain can take a time exponential in $D$ to mix, and since the posterior becomes spiked as $N$ grows, iterative algorithms can take a time exponential in $N$ to leave a neighborhood of a local optimum~\cite{nicklwang2024}. The preconditioned Crank Nicolson chain of Section~\ref{sec:dimrobust} has a spectral gap independent of the discretization~\cite{hairer2014}, but the bound is not quantitative in $N$ and $D$~\cite{nicklwang2024}. The question is when the posterior mean can be computed to accuracy $\varepsilon$ in a number of steps polynomial in $N$, $D$ and $\varepsilon^{-1}$, with high probability under the data. Nickl and Wang answer it for one model, under a hypothesis on the forward map called gradient stability. We ask where it holds, above all for the two dimensional Navier Stokes equations. Numbers of theorems, lemmas and equations of Nickl and Wang, of Konen and Nickl and of Castre and Nickl refer to the arXiv versions 2009.05298v2, 2507.18279v2 and 2601.10326v1.

\subsubsection{Setting} We use the notation of Nickl and Wang, with the forward map written $\mathcal{G}$, $P^X$ the uniform probability measure on a bounded set $\mathcal{X}$ and $\|\cdot\|_{L^2}$ the norm of $L^2(P^X)$. For $\theta\in\R^D$, the coefficients of the unknown in the first $D$ basis elements, the data $Z_i=(Y_i,X_i)$ satisfy
\begin{equation}\label{eq:gs_model}
Y_i=\mathcal{G}(\theta)(X_i)+\varepsilon_i,\qquad X_i\sim P^X,\quad\varepsilon_i\sim N(0,1),\qquad i=1,\dots,N,
\end{equation}
all independent, under a truth $\theta_0$ that need not lie in $\R^D$ (we write $\P_{\theta_0}^N$ and $\E_{\theta_0}$). For an evolution equation, $\mathcal{X}=[0,T]\times\T^d$ and $\mathcal{G}(\theta)(t,x)$ is the solution at time $t$ and position $x$. This is the observation model~\eqref{eq:observation_model} with $\mathcal{O}_j$ the evaluation at a random point of space and time and unit noise. With $\ell(\theta,Z)=-\frac12(Y-\mathcal{G}(\theta)(X))^2$, $\ell_N(\theta)=\sum_{i\le N}\ell(\theta,Z_i)$ and a Gaussian prior $N(0,\Sigma)$ with density $\pi$, the posterior $\Pi(\cdot\mid Z)$ has density proportional to $e^{\ell_N}\pi=e^{-H}$, where $H(\theta)=-\ell_N(\theta)+\frac12\theta^\top\Sigma^{-1}\theta$.

\subsubsection{Gradient stability} Differentiating $\ell$ twice and using $\E_{\theta_0}[Y\mid X]=\mathcal{G}(\theta_0)(X)$ gives, for every $v\in\R^D$,
\begin{equation}\label{eq:gs_hess}
\begin{aligned}
v^\top\E_{\theta_0}\big[-\nabla^2\ell(\theta,Z)\big]v={}&\big\|v^\top\nabla\mathcal{G}(\theta)\big\|_{L^2}^2\\
&+\big\langle\mathcal{G}(\theta)-\mathcal{G}(\theta_0),\,v^\top\nabla^2\mathcal{G}(\theta)v\big\rangle_{L^2}.
\end{aligned}
\end{equation}
At $\theta=\theta_0$, when this lies in $\R^D$, the second term vanishes and the matrix is the Fisher information $I(\theta_0)=\E_X[\nabla\mathcal{G}(\theta_0)(X)\nabla\mathcal{G}(\theta_0)(X)^\top]$ of one observation. The expected Hessian of $-\ell_N$ is $N$ times this matrix, so a lower bound on its smallest eigenvalue says that the data are informative in every direction, and by~\eqref{eq:gs_hess} it follows from an injectivity estimate for the gradient.

\begin{definition}[Gradient stability]\label{def:gs_stab}
Let $B\subset\R^D$ and $c>0$. The forward map $\mathcal{G}$ is gradient stable on $B$ with constant $c$ if $\theta\mapsto\mathcal{G}(\theta)(x)$ is differentiable on $B$ for every $x$ and
\begin{equation}\label{eq:gs_stab}
\big\|v^\top\nabla\mathcal{G}(\theta)\big\|_{L^2}\ge c\,|v|\qquad\text{for all }\theta\in B\text{ and }v\in\R^D.
\end{equation}
\end{definition}

Nickl and Wang take $B=\{\theta:|\theta-\theta^*|_1<\eta\}$, with $\theta^*\in\R^D$ a centering point (the projection of $\theta_0$) and $|\theta|_1^2=\theta^\top M\theta$ for a symmetric positive definite matrix $M$. The difficulty lies in the dependence of the constant $c$ on $D$, which Castre and Nickl~\cite{castrenickl2026} make explicit. What enters the theorem is the local curvature condition
\begin{equation}\label{eq:gs_curv}
\inf_{\theta\in B}\lambda_{\min}\big(\E_{\theta_0}[-\nabla^2\ell(\theta,Z)]\big)\ge c_{\min}>0 .
\end{equation}
It follows from~\eqref{eq:gs_stab} by a line that we add. If $k_0$ and $k_2$ bound the supremum norms of $\mathcal{G}(\theta)-\mathcal{G}(\theta_0)$ and of $\|\nabla^2\mathcal{G}(\theta)\|_{\mathrm{op}}$ over $B$, the second term of~\eqref{eq:gs_hess} is at least $-k_0k_2|v|^2$ by the Cauchy Schwarz inequality, so gradient stability with constant $c$ and $k_0k_2\le c^2/2$ give~\eqref{eq:gs_curv} with $c_{\min}=c^2/2$.

The Bernstein von Mises theory of this chapter rests on invertibility of the information operator, the derivative of the forward map composed with its adjoint, between Sobolev spaces, which Konen and Nickl~\cite{konennickl2026} prove for the two dimensional Navier Stokes equations. Gradient stability is its quantitative counterpart on $D$ coordinates. It is local, and hence weaker than the Polyak \L{}ojasiewicz condition of nonconvex optimization, a global lower bound on the squared gradient by the suboptimality, as Nickl and Wang remark.

\subsubsection{The theorem of Nickl and Wang} \emph{Local regularity} asks that $\mathcal{G}(\theta)(x)$ be twice continuously differentiable in $\theta\in B$, with $\mathcal{G}(\theta)-\mathcal{G}(\theta_0)$, $\nabla\mathcal{G}(\theta)$ and $\nabla^2\mathcal{G}(\theta)$ bounded in the supremum norm over $x$ by $k_0,k_1,k_2$ and Lipschitz in $\theta$ for $|\cdot|_1$ with constants $m_0,m_1,m_2$. \emph{Local curvature} asks for~\eqref{eq:gs_curv} and for a bound $c_{\max}\ge c_{\min}$ on the expected log likelihood and its first two derivatives over $B$. The algorithm needs an initializer $\theta_{\mathrm{init}}\in\R^D$ computed from the data and a surrogate $\tilde\ell_N=\alpha_\eta\ell_N-Kg_\eta$ of the log likelihood. Here $\alpha_\eta$ is a smooth cutoff in $|\theta-\theta_{\mathrm{init}}|_1/\eta$, $g_\eta$ is a convex penalty that vanishes for $|\theta-\theta_{\mathrm{init}}|_1\le\eta/2$, and $K>0$ is large. The surrogate posterior $\tilde\Pi(\cdot\mid Z)$ has density proportional to $e^{\tilde\ell_N}\pi$, and the unadjusted Langevin algorithm (ULA) for it, with step size $\gamma>0$ and independent $\xi_k\sim N(0,I_D)$, is
\begin{equation}\label{eq:gs_ula}
\vartheta_0=\theta_{\mathrm{init}},\qquad\vartheta_{k+1}=\vartheta_k+\gamma\big(\nabla\tilde\ell_N(\vartheta_k)-\Sigma^{-1}\vartheta_k\big)+\sqrt{2\gamma}\,\xi_{k+1}.
\end{equation}
Three events enter. On $E_{\mathrm{conv}}$, $\lambda_{\min}(-\nabla^2\ell_N(\theta))\ge Nc_{\min}/2$ and $|\ell_N|+|\nabla\ell_N|+\|\nabla^2\ell_N\|_{\mathrm{op}}\le N(5c_{\max}+1)$ for all $\theta\in B$. Further $E_{\mathrm{init}}=\{|\theta_{\mathrm{init}}-\theta^*|_1\le\eta/8\}$ and $E_{\mathrm{wass}}(\rho)=\{W_2^2(\Pi(\cdot\mid Z),\tilde\Pi(\cdot\mid Z))\le\rho/2\}$, where $W_2$ is the Wasserstein distance of Chapter~\ref{ch:entropy}.

\begin{theorem}[Nickl and Wang, in simplified form]\label{thm:gs_main}
Let the data follow~\eqref{eq:gs_model}, let the prior be $N(0,\Sigma)$ and let local regularity and local curvature hold on $B$ with $\eta\in(0,1]$. Let $R_N$ be the quantity of equation~(45) of~\cite{nicklwang2024}, a universal constant times $N$ times the minimum of eight ratios of the constants above; in particular $R_N\le CNc_{\min}^2/(k_1^2+k_0k_2+k_2)^2$.
\begin{enumerate}[label=(\roman*),nosep]
\item If $D\le R_N$, then $\P_{\theta_0}^N(E_{\mathrm{conv}})\ge1-32e^{-R_N}-e^{-N/8}$.
\item Let $D\le R_N$, $K\ge CN(c_{\max}+1)(1+\lambda_{\max}(M)\eta^{-2})/\lambda_{\min}(M)$, $m=\frac12Nc_{\min}+\lambda_{\min}(\Sigma^{-1})$ and $\Lambda=7K\lambda_{\max}(M)+\lambda_{\max}(\Sigma^{-1})$. Then for every $\rho>0$, $0<\gamma\le1/(m+\Lambda)$ and $k\ge0$, on $E=E_{\mathrm{conv}}\cap E_{\mathrm{init}}\cap E_{\mathrm{wass}}(\rho)$,
\begin{equation}\label{eq:gs_w2}
W_2^2\big(\Law(\vartheta_k),\Pi(\cdot\mid Z)\big)\le\rho+b(\gamma)+4\Big(\tau+\frac Dm\Big)\Big(1-\frac{\gamma m}2\Big)^k,
\end{equation}
where $b(\gamma)=c_1(\gamma D\Lambda^2m^{-2}+\gamma^2D\Lambda^4m^{-3})$, $\tau=c_2\kappa(\Sigma)(1+\eta^2/\lambda_{\min}(M)+R^2)$ for any $R\ge|\theta^*|$, $\kappa(\Sigma)=\lambda_{\max}(\Sigma)/\lambda_{\min}(\Sigma)$, and $C,c_1,c_2$ are universal constants.
\end{enumerate}
\end{theorem}

Part (i) is Lemma~3.4 and part (ii) is Proposition~3.6 with Theorem~3.7 of~\cite{nicklwang2024}. Their Theorem~3.7 allows $0<\gamma\le1/\Lambda$; its proof applies, through their Proposition~A.4, Theorem~5 of Durmus and Moulines~\cite{durmusmoulines2019}, which is proved for $\gamma\le1/(m+\Lambda)$, and we state (ii) in that range. Since $m\le\Lambda$, this costs at most a factor $2$ in the step, and for large $N$ the range contains the step sizes used below. We quote the theorem without proof. The theorem does not say that $E_{\mathrm{init}}$ and $E_{\mathrm{wass}}(\rho)$ have high probability. Bounding their probability, Nickl and Wang write, `may pose a formidable challenge in its own right', and they do it for the Schr\"odinger equation only.

\begin{corollary}[Posterior mean and mode]\label{cor:gs_mean}
In the setting of Theorem~\ref{thm:gs_main}(ii) let $\bar\vartheta$ be the average of $\vartheta_{J_{\mathrm{in}}+1},\dots,\vartheta_{J_{\mathrm{in}}+J}$.
\begin{enumerate}[label=(\alph*),nosep]
\item There are constants $c_3,c_5>0$ such that on $E$, for $t\ge\sqrt8\,(\rho+b(\gamma))^{1/2}$, $mJ\gamma\ge1$ and $J_{\mathrm{in}}\ge(c_3/m\gamma)\log\big(1+(\rho+b(\gamma))^{-1}+\tau+D/m\big)$, $\P_{\theta_{\mathrm{init}}}(|\bar\vartheta-\E^\Pi[\theta\mid Z]|\ge t)\le2D\exp(-c_5t^2m^2J\gamma/D)$.
\item If the posterior density has a maximizer and every maximizer lies in $\{|\theta-\theta^*|_1\le3\eta/8\}$, then on $E_{\mathrm{conv}}\cap E_{\mathrm{init}}$ the posterior has a unique mode, and gradient ascent on the logarithm of the surrogate density from $\theta_{\mathrm{init}}$ with step $(2\Lambda)^{-1}$ converges to it, with squared error at most a constant, which depends on the data, times $(1-m/(2\Lambda))^k$.
\end{enumerate}
\end{corollary}

Part (a) is Corollary~3.9 of~\cite{nicklwang2024}, quoted without proof, with the factor $1+1/(mJ\gamma)\le2$ in the exponent of their bound absorbed in $c_5$. Part (b) is our formulation of their Proposition~A.2 and of the argument after their equation~(158), and the argument is short. On $E_{\mathrm{conv}}\cap E_{\mathrm{init}}$ the surrogate $\tilde\ell_N$ and $\ell_N$ have the same gradient on the ball of (b), so a maximizer of the posterior density is a critical point of the logarithm of the surrogate density, which is strongly concave and has one critical point only. Proposition~A.2 is the geometric rate of gradient descent for a function whose Hessian lies between $mI$ and $\Lambda I$. With $\rho\le\varepsilon^2/16$ and $b(\gamma)\le\varepsilon^2/16$ the threshold for $t$ is at most $\varepsilon$, and the number of steps is polynomial in $N$, $D$, $m^{-1}$, $\Lambda$ and $\varepsilon^{-1}$.

\subsubsection{The Schr\"odinger equation} Nickl and Wang prove $E_{\mathrm{init}}$ and $E_{\mathrm{wass}}$ for the steady state Schr\"odinger equation $\frac12\Delta u-fu=0$ on a bounded smooth domain $\mathcal{O}\subset\R^d$, $d\le3$, of unit volume, with boundary values $u=g$, an elliptic boundary value problem. The unknown is $f_\theta=\Phi\circ F_\theta$, where $F_\theta=\sum_{k\le D}\theta_ke_k$ in the eigenfunctions of the Dirichlet Laplacian, with eigenvalues $\lambda_k$, and $\Phi$ is a regular link function. The forward map is $\mathcal{G}(\theta)=u_{f_\theta}$, the design is uniform on $\mathcal{O}$, and the prior is $N(0,N^{-d/(2\alpha+d)}\Lambda_\alpha^{-1})$ with $\Lambda_\alpha=\diag(\lambda_1^\alpha,\dots,\lambda_D^\alpha)$ and an integer $\alpha>6$.

\begin{theorem}[Nickl and Wang, Schr\"odinger equation, in simplified form]\label{thm:gs_schr}
Let $\theta_0$ satisfy $\sum_k\lambda_k^\alpha\theta_{0,k}^2<\infty$, and for some constants $c_0>0$ and $0<c_0'<1/2$ let $D\le c_0N^{d/(2\alpha+d)}$ and $\|\mathcal{G}(\theta_{0,D})-\mathcal{G}(\theta_0)\|_{L^2}\le c_0'N^{-\alpha/(2\alpha+d)}$, where $\theta_{0,D}$ is the projection of $\theta_0$ onto $\R^D$. Run~\eqref{eq:gs_ula} with $M=I$, $\theta^*=\theta_{0,D}$, $\eta=D^{-4/d}/\log N$, $K=ND^{8/d}(\log N)^3$, $\gamma\le(ND^{8/d}(\log N)^4)^{-1}$ and the initializer of Theorem~B.6 of~\cite{nicklwang2024}. On events of probability at least $1-ce^{-c'N^{d/(2\alpha+d)}}$, for every $P>0$ and $\varepsilon\ge N^{-P}$:
\begin{enumerate}[label=(\roman*),nosep]
\item there are a step size and a number of steps $O(N^{b_1}D^{b_2}\varepsilon^{-b_3})$ after which $W_2(\Law(\vartheta_k),\Pi(\cdot\mid Z))\le\varepsilon$, and averaging the chain computes the posterior mean to accuracy $\varepsilon$ at a cost of the same form; the exponents $b_1,b_2,b_3>0$ are not specified and may differ between the two statements;
\item the posterior has a unique mode, and gradient ascent on the logarithm of the surrogate density with step size $(ND^{8/d}(\log N)^4)^{-1}$ comes within $cN^{-\alpha^2/((2\alpha+d)(\alpha+2))}$ of $\theta_0$ in $\ell^2$ after $O(D^{12/d}(\log N)^5)$ iterations.
\end{enumerate}
\end{theorem}

We quote the theorem without proof. It collects Theorem~1.1 and Theorems~2.5, 2.7 and~2.8 of~\cite{nicklwang2024}, which are Theorem~\ref{thm:gs_main} and Corollary~\ref{cor:gs_mean} with the scalings of their Proposition~4.1, together with their Theorems~B.6 and~4.14 for $E_{\mathrm{init}}$ and $E_{\mathrm{wass}}$. The scalings are $k_0\simeq k_1\simeq1$, $k_2\simeq m_0\simeq m_1\simeq D^{2/d}$, $m_2\simeq D^{4/d}$, $c_{\max}\simeq1$ and, from an elliptic stability estimate, $c_{\min}\simeq D^{-4/d}$. Then $R_N\simeq ND^{-12/d}$, $m\simeq ND^{-4/d}$, $\Lambda\simeq ND^{8/d}(\log N)^3$ and $\Lambda/m\simeq D^{12/d}(\log N)^3$. Reading~\eqref{eq:gs_w2} with these values, accuracy $\tilde\varepsilon=\varepsilon m^{1/2}$ in units of the posterior scale takes about $D^{1+24/d}(\log N)^6\tilde\varepsilon^{-2}$ steps, of degree thirteen in $D$ for $d=2$.

\begin{remark}[Scope of the example]\label{rem:gs_scope}
The equation is elliptic and steady: there is no time, the unknown is a coefficient, and the design fills the domain. The guarantee is local and needs a warm start in a ball of radius of order $D^{-4/d}/\log N$ around $\theta_{0,D}$, which the initializer reaches by estimating $u$ by convex penalized least squares and inverting the equation, $f=\Delta\hat u/(2\hat u)$, at the price of two derivatives. A cold start can fail even for a unimodal target~\cite{bandeira2023}. Theorem~\ref{thm:gs_main} is therefore a template and not a theorem about every gradient stable forward map.
\end{remark}

\begin{proof}[Sketch of the proof of Theorem~\ref{thm:gs_main}]
\emph{(a) Local strong convexity of $-\ell_N$.} By Weyl's inequality it suffices to keep the empirical Hessian within $Nc_{\min}/2$ of its mean uniformly on $B$. A net of the unit sphere of mesh $1/4$, with at most $9^D$ points, reduces the operator norm to a maximum over the net, up to a factor two. For a fixed direction the centered quadratic form splits into its value at $\theta^*$, handled by Hoeffding's and Bernstein's inequalities, and an increment over $\theta\in B$, an empirical process that generic chaining controls (Lemma~3.12 of~\cite{nicklwang2024}). The probability bounds beat the union bound $9^D$ when $D\le R_N$. This is part~(i).

\emph{(b) A strongly log concave surrogate.} On $E_{\mathrm{conv}}\cap E_{\mathrm{init}}$, $\tilde\ell_N=\ell_N$ on $\{|\theta-\theta^*|_1\le3\eta/8\}$, $-\nabla^2\tilde\ell_N\succeq\frac12Nc_{\min}I$ on all of $\R^D$ and $\nabla\tilde\ell_N$ is $7K\lambda_{\max}(M)$ Lipschitz (Proposition~3.6 of~\cite{nicklwang2024}). Within distance $3\eta/4$ of $\theta_{\mathrm{init}}$ the cutoff equals one and the penalty is convex, so $-\nabla^2\tilde\ell_N$ inherits the curvature that $E_{\mathrm{conv}}$ gives. Beyond it, on $E_{\mathrm{conv}}$ the Hessian of $\alpha_\eta\ell_N$ is at most of order $N(c_{\max}+1)(1+\lambda_{\max}(M)\eta^{-2})$ in norm, while the Hessian of $Kg_\eta$ is at least $K\lambda_{\min}(M)/3$, so the condition on $K$ lets the penalty dominate. With the prior, $\tilde U=-\tilde\ell_N+\frac12\theta^\top\Sigma^{-1}\theta$ satisfies $mI\preceq\nabla^2\tilde U\preceq\Lambda I$.

\emph{(c) Langevin sampling of a strongly log concave target.} Let $\mu\propto e^{-U}$ with $mI\preceq\nabla^2U\preceq\Lambda I$. The law of the Langevin diffusion $\dd L_t=-\nabla U(L_t)\dd t+\sqrt2\,\dd W_t$ solves the Fokker Planck equation~\eqref{eq:fp} with $V=U$, and by Theorem~\ref{thm:jko} it is the Wasserstein gradient flow of $\KL(\cdot\,\|\,\mu)$. The criterion of Bakry and \'Emery~\cite{bakry1985}, stated after Theorem~\ref{thm:lsi}, turns $\nabla^2U\succeq mI$ into a log Sobolev inequality with constant $2/m$, hence $\KL(\rho_t\|\mu)\le e^{-2mt}\KL(\rho_0\|\mu)$, and Talagrand's inequality $W_2^2\le(2/m)\KL$ transfers the decay to $W_2$. The discrete time bound that~\eqref{eq:gs_ula} needs is due to Durmus and Moulines~\cite{durmusmoulines2019}, after~\cite{dalalyan2017,durmusmoulines2017}. The chain starts at a point, whose relative entropy is infinite, so the bound is in $W_2$: for $0<\gamma\le1/(m+\Lambda)$ and $\theta_U$ the minimizer of $U$ (Theorem~5 of~\cite{durmusmoulines2019}, in the form of Proposition~A.4 of~\cite{nicklwang2024}),
\[
\begin{aligned}
W_2^2(\Law(\vartheta_k),\mu)\le{}&2\Big(1-\frac{m\gamma}2\Big)^k\Big(|\theta_{\mathrm{init}}-\theta_U|^2+\frac Dm\Big)\\
&+\frac{18\gamma D\Lambda^2}{m^2}+\frac{6\gamma^2D\Lambda^4}{m^3}.
\end{aligned}
\]
The first term decays geometrically in $k$, as $e^{-2mt}$ does in $t$ for the diffusion, but at the weaker rate $m\gamma/2$ per step instead of $2m\gamma$. The others, which have no counterpart in Chapter~\ref{ch:entropy}, are the bias of the Euler scheme. With $\gamma\simeq\varepsilon^2m^2/(D\Lambda^2)$ the bias is of order $\varepsilon^2$, and $k\simeq(\Lambda/m)^2D(\varepsilon^2m)^{-1}\log(\cdot)$ steps suffice.

\emph{(d) From the surrogate to the posterior.} The penalty puts the maximizer of $\tilde\ell_N$ within a distance of order $1+\eta\lambda_{\min}(M)^{-1/2}$ of $\theta_{\mathrm{init}}$, and the prior moves the mode of $\tilde\Pi$ by an amount controlled by $\kappa(\Sigma)$, so $|\theta_{\mathrm{init}}-\theta_U|^2\lesssim\tau$ for $\mu=\tilde\Pi$. The triangle inequality for $W_2$ and $W_2^2(\Pi,\tilde\Pi)\le\rho/2$ then give~\eqref{eq:gs_w2}. For the Schr\"odinger equation Theorem~4.14 of~\cite{nicklwang2024} bounds $E_{\mathrm{wass}}$: $W_2^2$ is at most twice the integral of $|\theta-\theta_1|^2$ against $|\Pi-\tilde\Pi|$ for every point $\theta_1$ (Theorem~6.15 of~\cite{villani2009}), and for $\theta_1$ the posterior mode the two measures agree up to normalizing constants near $\theta_1$ and carry exponentially small mass away from it.
\end{proof}

\subsubsection{Where the log Sobolev inequality enters} The proof uses the theory of Chapter~\ref{ch:entropy} in step (c) and only there, and for the surrogate and not for the posterior, which is not strongly log concave and for which Nickl and Wang prove no log Sobolev inequality. Results that assume a log Sobolev or Poincar\'e inequality for the target~\cite{chewi2025}, or that it is convex outside a ball~\cite{ma2019}, yield bounds exponential in $N$ when applied to the posterior without further work, as Nickl and Wang observe: the Holley Stroock principle~\cite{holleystroock1987} costs a factor exponential in the oscillation of $\ell_N$, which is of order $N$ for a bounded forward map, and the bound of Ma and coauthors is exponential in the Lipschitz constant of $\nabla\ell_N$. The theorem rests on strong log concavity of the surrogate, its closeness to the posterior and the initializer.

\subsubsection{Where gradient stability is known} Nickl and Wang verify gradient stability for the steady state Schr\"odinger equation (their Lemma~4.7), and Bohr and Nickl~\cite{bohrnickl2024} extend the theory and verify its stability conditions for the non Abelian X ray transform. Castre and Nickl~\cite{castrenickl2026} list further parameter identification problems with linear equations where gradients of forward maps have been computed, among them Darcy's problem~\cite{nickl2023book,altmeyer2022}, and initial value to solution maps of reaction diffusion and Navier Stokes equations, see the next subsection. They write that these proofs do not extend straightforwardly to a general parameter of a nonlinear equation, where two nonlinearities meet. Their device is the implicit function theorem: $f(\theta,\mathcal{G}(\theta))=0$ gives $D\mathcal{G}=-(D_2f)^{-1}D_1f$, and gradient stability reduces to injectivity of $D_1f$. For a reaction diffusion equation with unknown reaction term they prove differentiability of the parameter to solution map (their Theorem~2.1) and only sketch gradient stability (Remark~2.2).

\begin{remark}[Castre and Nickl: the McKean Vlasov equation]\label{rem:gs_cn}
Let $\rho_W$ solve $\partial_t\rho=\Delta\rho+\nabla\cdot(\rho\nabla W*\rho)$ on $[0,T]\times\T^d$ with a smooth initial condition $\varphi$ bounded below and an unknown interaction potential $W$ of degree at most $K$ in the trigonometric basis, so $D\simeq K^d$, and let $\rho_W$ be observed as in~\eqref{eq:gs_model}. Assume that $|\hat\varphi_k|\ge c_*|k|^{-\zeta}$ for a large enough $\zeta$, which excludes the steady state $\varphi=1$~\cite{nickl2025}. Gradient stability holds with constant of order $K^{-3\zeta}$ (their Theorem~2.14), and $c_{\min}\ge c_0D^{-6\zeta/d}$ on a ball around the truth of radius polynomially small in $D$ (Theorem~3.3). The posterior is $W_2$ close to a log concave surrogate with probability tending to one (Theorem~3.10), and \emph{if a feasible initializer exists} ULA mixes in a time polynomial in $\varepsilon^{-1}$, $D$ and $N$ (Theorem~3.11). The initializer is assumed, not constructed, and the unknown is a parameter and not an initial condition.
\end{remark}

\subsubsection{The initial condition problem} For the initial condition of the two dimensional Navier Stokes equations, Nickl and Titi~\cite{nickltiti2024} prove posterior consistency at inverse logarithmic rates, and show that no faster rate is possible when the observation times lie in $[T_0,T]$ with $T_0>0$. For initial conditions on the global attractor of a two dimensional periodic reaction diffusion system, Konen and Nickl~\cite{konennickl2026attractor} show that the initial condition and the states can be recovered from discrete measurements at near parametric rates. For times sampled from $[0,T]$, Nickl~\cite{nicklbvm2024} proves a Bernstein von Mises theorem for the trajectory of dissipative systems such as reaction diffusion, and Konen and Nickl~\cite{konennickl2026} prove one for Navier Stokes, with polynomial contraction rates for the initial condition (their Theorem~2.2) and invertibility of the information operator (Theorem~2.6). None of this is a theorem about computation. Konen and Nickl say only that provably reliable computation `can be expected to be possible in polynomial time in relevant parameters' by an MCMC algorithm combining the ideas of Nickl and Wang with their estimates, started at the boundary trace of a smoothing estimator, and that the details will be studied elsewhere. We have not found a later paper that turns these expectations into theorems.

What the estimates of Konen and Nickl do give is the first ingredient. The velocity is vector valued, and Remark~3.11 of~\cite{nicklwang2024} extends the theory to values in a finite dimensional inner product space, with the norm of that space in the likelihood and in $\|\cdot\|_{L^2}$. The following statement is our deduction from their Propositions~3.10 and~3.12, and not a result of their paper.

\begin{proposition}[Gradient stability for the initial condition of two dimensional Navier Stokes, with time zero observed]\label{prop:gs_ns}
Let $\nu,T,L>0$ and $f\in\Hds1$, let $\mathcal{X}=[0,T]\times\T^2$ with the uniform probability measure, let $e_1,e_2,\dots$ be the eigenfunctions of the Stokes operator on $\T^2$ with eigenvalues $\lambda_1\le\lambda_2\le\cdots$, and let $\mathcal{G}(\theta)(t,x)=u_\theta(t,x)$ be the velocity of the Navier Stokes solution with viscosity $\nu$, forcing $f$ and initial condition $\sum_{k\le D}\theta_ke_k$. There is $c>0$, depending only on $\nu$, $T$, $\|f\|_{\Hds1}$, $L$ and the torus, such that $\|v^\top\nabla\mathcal{G}(\theta)\|_{L^2(\mathcal{X})}\ge c\,D^{-1/2}|v|$ for every $D$, every $v\in\R^D$ and every $\theta\in\R^D$ with $\|\sum_k\theta_ke_k\|_{\Hds2}\le L$.
\end{proposition}

\begin{proof}
By Proposition~3.12 of Konen and Nickl, $\theta\mapsto u_\theta$ is differentiable with derivative the linearization $\I_\theta$ of the flow, so $v^\top\nabla\mathcal{G}(\theta)(t,\cdot)=\I_\theta[\xi_v](t)$ for $\xi_v=\sum_{k\le D}v_ke_k$. Their Proposition~3.10, with $a=-1$ and $a^*=2$, gives $\|\xi\|_{\Hds{-1}}^2\le C\int_0^T\|\I_\theta[\xi](t)\|_{L^2}^2\dd t$ for $\|\theta\|_{\Hds2}\le L$, with $C$ depending on $\nu,T,\|f\|_{\Hds1},L$. Hence, with the uniform probability measure on $\mathcal{X}$, $\|v^\top\nabla\mathcal{G}(\theta)\|_{L^2(\mathcal{X})}^2=(T|\T^2|)^{-1}\int_0^T\|\I_\theta[\xi_v](t)\|_{L^2}^2\dd t\ge(CT|\T^2|)^{-1}\|\xi_v\|_{\Hds{-1}}^2=(CT|\T^2|)^{-1}\sum_kv_k^2/\lambda_k\ge(CT|\T^2|)^{-1}|v|^2/\lambda_D$, and $\lambda_D\asymp D$ by Weyl's law in two dimensions.
\end{proof}

The constant is polynomial in $D$, which is what a design that includes time zero buys, but it is only the first ingredient. Curvature on a ball needs $k_0k_2\le c^2/(2D)$ there, and local regularity needs derivatives of the flow up to order three in the first $D$ coordinates, with bounds polynomial in $D$. Konen and Nickl bound the flow and its first derivative in Sobolev norms (their Propositions~3.7 and~3.9), which with the Sobolev embedding of $\Hds2$ into the continuous functions gives $k_0$, $k_1$ and $m_0$ polynomial in $D$ for $\theta$ in a ball of $\Hds2$, because $\|\sum_{k\le D}v_ke_k\|_{\Hds2}\lesssim\lambda_D|v|$, and they bound the quadratic Taylor remainder (their Proposition~3.12). We have not found bounds for the second and third derivatives of the flow, which $k_2$, $m_1$ and $m_2$ need. The warm start needs an initializer within the radius $\eta/8$ of the projection of the truth, and the boundary trace of a smoothing estimator is a candidate whose accuracy at time zero has to be proved. The event $E_{\mathrm{wass}}$ needs exponentially small posterior mass outside the ball, whereas the contraction theorem of Konen and Nickl holds in probability. Whether these steps can be taken is open as far as we know.

\subsubsection{A computation for the heat equation} Observation that begins after time zero changes the constant. Take $\mathcal{G}(\theta)(t,\cdot)=\sum_{k\le D}e^{-k^2t}\theta_k\phi_k$, the heat equation of Section~\ref{sec:linearheat} with the basis orthonormal in $L^2(P^X)$. It is linear, so $\nabla^2\mathcal{G}=0$ and by~\eqref{eq:gs_hess} the expected Hessian of $-\ell_N$ is $N\diag(\iota_1,\dots,\iota_D)$ at every $\theta$, where $\iota_k=\E[e^{-2k^2t}]$ is the information that one observation carries about mode $k$. Gradient stability holds on all of $\R^D$ with $c^2=c_{\min}=\min_{k\le D}\iota_k$. At a single time $T$, $\iota_k=e^{-2k^2T}$. For times uniform on $[T_0,T]$ it is the quantity of~\eqref{m:eq:iota} with $\nu\lambda_j=k^2$, of order $(2k^2T)^{-1}$ if $T_0=0$ and of order $e^{-2k^2T_0}$, up to a factor polynomial in $k$, if $T_0>0$. For $T=0.05$ and $k=40$ it is $6.25\times10^{-3}$, $3.5\times10^{-20}$ and $2.3\times10^{-37}$ for $T_0=0$, $0.0125$ and $0.025$. Only the first is polynomial in $D$.

The consequence for Theorem~\ref{thm:gs_main} is a restriction on $D$. At a single time $T$, $k_2=0$ and $k_1$ is bounded, so $R_N\le CNc_{\min}^2/k_1^4$ and $D\le R_N$ forces $De^{4D^2T}\lesssim N$, that is $D\lesssim\sqrt{\log N/(4T)}$ up to lower order terms, and $D\lesssim(\log N)^{d/2}$ in $d$ dimensions. For $N$ of order $\delta^{-2}$ this is $K_\delta/\sqrt2$, where $K_\delta=\sqrt{\log(1/\delta)/T}$ is the mode up to which the data dominate the prior in the proof of Theorem~\ref{thm:heat_rates}. For $T=0.05$ and $\delta=0.002$, $K_\delta\approx11.1$ and $K_\delta/\sqrt2\approx7.9$, up to the unspecified constant $C$. A gradient stability constant that is exponentially small in $D^{2/d}$ therefore lets the theorem apply only up to about the dimension that the data resolve in any case, and the curvature beyond it comes from the prior, through $\lambda_{\min}(\Sigma^{-1})$ in $m$.

\subsubsection{Langevin and generators} Chapter~\ref{ch:generators} and Part~\ref{part:manifold} take another route to the same posterior: a generator, trained on simulations from the prior, maps noise to posterior draws (Table~\ref{tab:gs_compare}).

\begin{table}[t]\centering\footnotesize\setlength{\tabcolsep}{3pt}
\caption{Two routes to the posterior of an initial condition. $M$ is the number of simulated pairs and $J$ the dimension of the summary of Chapter~\ref{ch:adjoint}.}\label{tab:gs_compare}
\begin{tabular}{@{}>{\raggedright\arraybackslash}p{0.72in}>{\raggedright\arraybackslash}p{0.84in}>{\raggedright\arraybackslash}p{0.84in}>{\raggedright\arraybackslash}p{0.80in}>{\raggedright\arraybackslash}p{0.90in}@{}}
\toprule
Method & Guarantee & Assumed & Cost & Navier Stokes\\
\midrule
Warm start ULA, Theorem~\ref{thm:gs_main} & $W_2$ error $\varepsilon$ in polynomial time, with high probability & Gradient stability and local regularity, warm start, $E_{\mathrm{wass}}$ & $k$ forward and adjoint solves; nothing reused & Gradient stability if $t=0$ is sampled; the rest is open\\
\addlinespace
Gaussian head, Theorem \ref{m:thm:qnnA} & $\sqrt N$ times the $W_1$ error tends to zero & Assumptions \ref{m:ass:U} (open) and \ref{m:ass:u2p} (proved in the setting of Theorem~\ref{m:thm:u2ploc}) and a bound on the summary error; $M\gg J$ & $M$ solves and one fit, paid once; then a summary and one pass per draw & Conditional; six mode experiment (Section~\ref{m:sec:expns})\\
\bottomrule
\end{tabular}
\end{table}

In the heat equation experiment of Table~\ref{m:tab:cost} a generator used $8000$ solves in total and a pCN chain $2\times10^5$ solves for each data set. The Langevin guarantee concerns the data set in hand, in $W_2$ at any accuracy $\varepsilon\ge N^{-P}$, while a generator with a Gaussian head is a Monte Carlo estimate of the Laplace approximation and cannot beat it. For Langevin sampling, $E_{\mathrm{wass}}(\rho)$ asks that the posterior be $W_2$ close to a log concave measure, so a posterior with several well separated modes of comparable mass is outside the reach of the theorem. We do not know whether the posterior of a flow with filamenting vortices has several modes: in the computed example of Section~\ref{sec:simns2d} (Figure~\ref{fig:ns2d_posterior}), a flow smaller than that of Figure~\ref{fig:ns2d}, four chains started from independent prior draws agree. A Gaussian head cannot represent several modes either, and for a general generator Theorem~\ref{m:thm:qnnB} gives no uniform guarantee.

\subsubsection{What remains open} Most of the route from the theorem to the Navier Stokes equations is open. The steps listed after Proposition~\ref{prop:gs_ns} are not available, as far as we know, nor is a polynomial time guarantee for observation times bounded away from zero beyond dimensions of order $(\log N)^{d/2}$, or for posteriors with several modes. A polynomial time guarantee for fixed sensors and for priors other than Gaussian is the companion of Open Problem~\ref{prob:inference}.
